%LTeX: language=fr-FR
\documentclass[12pt]{article}
\usepackage[margin=0.8in]{geometry}
\usepackage[utf8]{inputenc}
\usepackage{graphicx}
\usepackage[export]{adjustbox}
\usepackage{amsmath}
\usepackage{amsfonts}
\usepackage{amssymb}
\usepackage[version=4]{mhchem}
\usepackage{stmaryrd}
\usepackage{bbold}
%\usepackage[fallback]{xeCJK}
%\usepackage{polyglossia}
\usepackage{fontspec}
\usepackage{censor}
\usepackage{currfile}
%\setCJKmainfont{Noto Serif CJK TC}
\usepackage{enumitem}
\usepackage{tikz}
\usetikzlibrary {arrows.meta, quotes, calc} 
\usepackage{pgfplots}
\pgfplotsset{compat=1.18}
\usepackage{amsthm}
\usepackage{tcolorbox}
\usepackage{array}
\usepackage{tabularray}
\usepackage{makecell}
\usepackage{esvect}
%\usepackage{changepage}   % for the adjustwidth environment
\usepackage{enumitem}
\usepackage{tasks}
\usepackage{fancyhdr}
\usepackage[french]{babel}
\usepackage{lastpage}
\usepackage{needspace}
\usepackage{tabularx}

\tcbuselibrary{theorems}

\definecolor{GrayCen}{HTML}{d2d2d2}
\definecolor{GrayBG}{HTML}{d9d9d9}
\definecolor{BlueDef}{HTML}{104e8b}
\definecolor{GreenProp}{HTML}{025540}
\definecolor{RedMet}{HTML}{cd5556}
\definecolor{BlackSav}{HTML}{000000}
\def\censorcolor{GrayCen}
\let\svcensorrule\censorrule
\renewcommand\censorrule[1]{%
\textcolor{\censorcolor}{\svcensorrule{#1}}}

\newtcbtheorem
[]% init options
{definition}% name
{Définition}% title
{%
  colback=GrayBG,
  colframe=BlueDef,
  fonttitle=\bfseries,
}% options
{def}% prefix

\newtcbtheorem
[]% init options
{proposition}% name
{Proposition}% title
{%
  colback=GrayBG,
  colframe=GreenProp,
  fonttitle=\bfseries,
}% options
{prop}% prefix

\newtcbtheorem
[]% init options
{method}% name
{Méthode}% title
{%
  theorem name,
  colback=GrayBG,
  colframe=RedMet,
  fonttitle=\bfseries,
}% options
{met}% prefix

\newtcbtheorem
[]% init options
{savoir}% name
{Savoir-faire du chapitre}% title
{%
  theorem name,
  colback=GrayBG,
  colframe=BlackSav,
  fonttitle=\bfseries,
}% options
{sav}% prefix

\newtheorem*{exampleil}{Exemple}

\newtheoremstyle{break}
{\topsep}{\topsep}%
{}{}%
{\bfseries}{}%
{\newline}{}%
\theoremstyle{break}
\newtheorem*{example}{Exemple.}
\newtheorem*{examples}{Exemples.}
\newtheorem*{remark}{Remarque.}
\newtheorem*{remarks}{Remarques.}

\newenvironment{demo}[1][\proofname]{%
  \begin{proof}[#1]$ $\newline
  }{%
  \end{proof}
}

%\setmainfont{CMU Serif}name

\title{
  \vspace{4em}
  \centering
  \tcbset{colframe=BlackSav,colback=GrayBG}
  \tcbox[tikznode]{
    Chapitre 4 \\Vecteurs dans un repère
  }
}

\author{}
\date{}

\pagestyle{empty}
\graphicspath{ {../Images/} }
\input{\currfiledir ../def.tex}
\def \Document {Cours}


\setlength\parindent{0pt}
%\setlist[itemize]{topsep=0pt}
%\setlist{nolistsep}
%\setlist[itemize]{noitemsep}
%\setlist[enumerate]{noitemsep}
\setlist[itemize,1]{label=$\bullet$}

\tikzset{cross/.style={path picture={
  \draw[black]
    (path picture bounding box.south east)--(path picture bounding box.north west)
    (path picture bounding box.south west)--(path picture bounding box.north east);
}}}

\newcolumntype{C}[1]{>{\centering\arraybackslash}m{#1}}
\newcolumntype{L}[1]{>{\raggedright\arraybackslash}m{#1}}
\newcolumntype{N}{@{}m{0pt}@{}}
\newcolumntype{Y}{>{\centering\arraybackslash}X}

\fancypagestyle{firststyle}
{
  \fancyhf{}
  \fancyhead[L]{\Structure}
  \fancyhead[R]{\Niveau}
  \fancyfoot[L]{Alexandre Fernandez}
  \fancyfoot[R]{{\thepage\  / \pageref{LastPage}}}
   %\renewcommand{\headrulewidth}{0pt} % removes horizontal header line
}


\begin{document}

\maketitle
\thispagestyle{firststyle}

\tableofcontents

% On crée une boîte temporaire pour stocker le texte
\newsavebox{\hidebox}

% -----------
% Macros des trous : version élève si \ELEVE est défini à la compilation
% (recette « prof + élève »), sinon version prof.
\ifdefined\ELEVE
% Version élève : contenu invisible
\newcommand{\hideM}[1]{\mbox{\vphantom{$#1$}\smash{\colorbox{gray!50}{\phantom{$\displaystyle #1$}}}} }
\newcommand{\hideMW}[1]{\mbox{\vphantom{$#1$}\smash{\colorbox{white}{\phantom{$\displaystyle #1$}}}} }
\newcommand{\hide}[1]{\mbox{\vphantom{#1}\smash{\colorbox{gray!50}{\phantom{#1}}}}}
\newcommand{\hideW}[1]{\mbox{\vphantom{#1}\smash{\colorbox{white}{\phantom{#1}}}}}
\NewEnviron{multihide}{%
  \colorbox{gray!50}{
    \begin{lrbox}{\hidebox}%
        \begin{minipage}{\linewidth}%
          \BODY
        \end{minipage}%
    \end{lrbox}%
    % On affiche un espace invisible ayant les dimensions exactes de la boîte
    \vrule width 0pt height \ht\hidebox depth \dp\hidebox\hskip\wd\hidebox%
  }
}
\NewEnviron{multihideW}{%
  \colorbox{white}{
    \begin{lrbox}{\hidebox}%
        \begin{minipage}{\linewidth}%
          \BODY
        \end{minipage}%
    \end{lrbox}%
    % On affiche un espace invisible ayant les dimensions exactes de la boîte
    \vrule width 0pt height \ht\hidebox depth \dp\hidebox\hskip\wd\hidebox%
  }
}
\newcommand{\hidetk}[1]{}
\else
% Version prof : contenu visible sur fond gris
\newcommand{\hideM}[1]{\mbox{\vphantom{$#1$}\smash{\colorbox{gray!50}{$\displaystyle #1$}}} }
\newcommand{\hideMW}[1]{\mbox{\vphantom{$#1$}\smash{\colorbox{white}{$\displaystyle #1$}}} }
\newcommand{\hide}[1]{\mbox{\vphantom{#1}\smash{\colorbox{gray!50}{#1}}}}
\newcommand{\hideW}[1]{\mbox{\vphantom{#1}\smash{#1}}}
\NewEnviron{multihide}{%
  \colorbox{gray!50}{
    \begin{minipage}{\linewidth}%
      \BODY
    \end{minipage}%
  }
}
\NewEnviron{multihideW}{%
  \colorbox{white}{
    \begin{minipage}{\linewidth}%
      \BODY
    \end{minipage}%
  }
}
\newcommand{\hidetk}[1]{#1}
\fi
% -----------

\newcommand{\wvec}[1]{\overrightarrow{#1}}
\newcommand{\ssi}[0]{\emph{ssi}\ }

% Macros pour la représentation des intervalles sur la droite réelle
% (crochet gauche fermé / ouvert, crochet droit fermé / ouvert)
\newcommand{\droite}{\draw[<->, >=Stealth] (0,0) -- (6,0);}
\newcommand{\brGF}[1]{\node[inner sep=0pt] at (#1,0) {\Large$[$};}
\newcommand{\brGO}[1]{\node[inner sep=0pt] at (#1,0) {\Large$]$};}
\newcommand{\brDF}[1]{\node[inner sep=0pt] at (#1,0) {\Large$]$};}
\newcommand{\brDO}[1]{\node[inner sep=0pt] at (#1,0) {\Large$[$};}
\newcommand{\etiq}[2]{\node[below=10pt,inner sep=0pt] at (#1,0) {\scriptsize$#2$};}


% ---------- Graphiques du cours ----------
% Quadrillage, axes et graduations : \axesG{xmin}{xmax}{ymin}{ymax}
\newcommand{\axesG}[4]{%
  \draw[gray!35, very thin, step=1] (#1,#3) grid (#2,#4);
  \draw[-{Latex[length=2mm]}] (#1,0) -- (#2+0.3,0);
  \draw[-{Latex[length=2mm]}] (0,#3) -- (0,#4+0.3);
  \foreach \x in {#1,...,#2} {\ifnum\x=0 \else \node[below, fill=white, fill opacity=0.7, text opacity=1, inner sep=1pt] at (\x,-0.18) {\small $\x$};\fi}
  \foreach \y in {#3,...,#4} {\ifnum\y=0 \else \node[left, fill=white, fill opacity=0.7, text opacity=1, inner sep=1pt] at (-0.1,\y) {\small $\y$};\fi}
  \node[below left, inner sep=1pt] at (-0.05,-0.18) {\small $0$};
}
\newcommand{\croix}[1]{\node[inner sep=0pt] at (#1) {$\times$};}

% Texte d'une étape dans les tableaux à deux colonnes : aligné à gauche et en haut
\newcommand{\txt}[1]{\multicolumn{1}{@{}p{0.48\linewidth}@{}}{\raggedright #1}}

% Cadre de réponse : blanc et vide dans la version élève (même hauteur)
\NewEnviron{reponse}{%
  \par\smallskip\noindent
  \begin{lrbox}{\hidebox}\begin{minipage}{\dimexpr\linewidth-2\fboxsep-2\fboxrule\relax}\BODY\end{minipage}\end{lrbox}%
  \fbox{\ifdefined\ELEVE\phantom{\usebox{\hidebox}}\else\usebox{\hidebox}\fi}%
  \par\smallskip
}

% ---------- Vecteurs ----------
% Coordonnées d'un vecteur en colonne
\newcommand{\col}[2]{{\renewcommand{\arraystretch}{1}\begin{pmatrix} #1 \\ #2 \end{pmatrix}}}
\newcommand{\colvide}{\col{\dots}{\dots}}
% Flèche de vecteur et point nommé
\tikzset{vec/.style={-{Stealth[length=3mm]}, very thick}}
\newcommand{\pt}[3]{\fill (#1) circle (2pt); \node[#2] at (#1) {$#3$};}


\newpage

\pagestyle{fancy}

\fancyhf{} % clear existing header/footer entries
\fancyhead[R]{ {\Niveau}}
\fancyhead[L]{ {\Chapitre}}
\fancyfoot[L]{ A. Fernandez }
\fancyfoot[R]{{\thepage\  / \pageref{LastPage}}}






\section{Repère du plan}

\begin{definition}{Repère}{}
  \begin{multihide}
  Soit $O$ un point du plan et $\vec\imath$, $\vec\jmath$ deux vecteurs de directions différentes.
  Le triplet $(O\,;\,\vec\imath,\vec\jmath)$ est un \textbf{repère} du plan, et le couple $(\vec\imath,\vec\jmath)$ est sa \textbf{base}.
  \begin{itemize}
    \item Si $\vec\imath$ et $\vec\jmath$ ont des directions perpendiculaires, le repère est \textbf{orthogonal}.
    \item Si $\vec\imath$ et $\vec\jmath$ ont la même norme, le repère est \textbf{normé}.
    \item Si les deux conditions sont vérifiées, le repère est \textbf{orthonormé}.
    \item Sinon, le repère est \textbf{quelconque}.
  \end{itemize}
  \end{multihide}
\end{definition}

\begin{example}
  Donner la nature de chaque repère.

  \begin{center}
  \begin{tabular}{ccc}
    \begin{tikzpicture}[scale=1.2]
      \draw[gray!35, very thin, xstep=1, ystep=0.5] (-1,-1) grid (3,1.5);
      \draw[vec] (0,0) -- (1,0) node[midway, below] {$\vec\imath$};
      \draw[vec] (0,0) -- (0,0.5) node[midway, left] {$\vec\jmath$};
      \node[below left] at (0,0) {$O$};
    \end{tikzpicture}
    &
    \begin{tikzpicture}[scale=1.2]
      \begin{scope}
        \clip (-1.5,-1) rectangle (2.5,1.5);
        \foreach \k in {-4,...,4} \draw[gray!35, very thin] ({\k-2*0.6},-2*0.8) -- ({\k+2*0.6},2*0.8);
        \foreach \k in {-2,...,2} \draw[gray!35, very thin] ({-3+\k*0.6},\k*0.8) -- ({4+\k*0.6},\k*0.8);
      \end{scope}
      \draw[vec] (0,0) -- (1,0) node[midway, below] {$\vec\imath$};
      \draw[vec] (0,0) -- (0.6,0.8) node[midway, left] {$\vec\jmath$};
      \node[below left] at (0,0) {$O$};
    \end{tikzpicture}
    &
    \begin{tikzpicture}[scale=1.2]
      \draw[gray!35, very thin, step=1] (-1,-1) grid (2,1.5);
      \draw[vec] (0,0) -- (1,0) node[midway, below] {$\vec\imath$};
      \draw[vec] (0,0) -- (0,1) node[midway, left] {$\vec\jmath$};
      \node[below left] at (0,0) {$O$};
    \end{tikzpicture}
    \\[0.5em]
    Repère \hide{orthogonal} & Repère \hide{quelconque} & Repère \hide{orthonormé}
  \end{tabular}
  \end{center}
\end{example}

\begin{remark}
  Dans toute la suite du chapitre, on se place dans un \textbf{repère orthonormé} $(O\,;\,\vec\imath,\vec\jmath)$.
\end{remark}

\section{Coordonnées d'un vecteur}

\subsection{Définition}

\begin{definition}{Coordonnées d'un vecteur}{}
  \begin{multihide}
  Pour tout vecteur $\vec u$ du plan, il existe deux réels \textbf{uniques} $x$ et $y$ tels que $\vec u = x\,\vec\imath + y\,\vec\jmath$.

  Ces réels sont les \textbf{coordonnées} de $\vec u$. On note $\vec u\col{x}{y}$, en colonne.
  \end{multihide}
\end{definition}

\begin{remarks}
  \begin{itemize}
    \item[]
    \item $x$ est le déplacement \textbf{horizontal} (négatif vers la gauche) et $y$ le déplacement \textbf{vertical} (négatif vers le bas).
    \item On écrit les coordonnées d'un vecteur \textbf{en colonne}, et celles d'un point \textbf{en ligne} : $A(x\,;\,y)$.
  \end{itemize}
\end{remarks}

\Needspace{8\baselineskip}
\begin{example}
  \begin{minipage}{0.45\linewidth}
    Lire les coordonnées des vecteurs représentés ci-contre.

    \medskip
    $\vec u\,\hideM{\col{4}{-2}}$ \qquad $\vec v\,\hideM{\col{-2}{-4}}$ \qquad $\vec w\,\hideM{\col{0}{-3}}$

    \medskip
    Représenter le vecteur $\vec t\col{3}{1}$ à partir du point $D$.
  \end{minipage}
  \hfill
  \begin{minipage}{0.52\linewidth}
    \begin{center}
    \begin{tikzpicture}[x=0.65cm, y=0.65cm]
      \axesG{-4}{4}{-4}{4}
      \draw[vec] (-3,4) -- (1,2) node[midway, above right] {$\vec u$};
      \draw[vec] (3,2) -- (1,-2) node[midway, right] {$\vec v$};
      \draw[vec] (-2,-1) -- (-2,-4) node[midway, left] {$\vec w$};
      \pt{-3,1}{left}{D}
      \hidetk{\draw[vec, blue!70!black] (-3,1) -- (0,2) node[midway, above left] {$\vec t$};}
    \end{tikzpicture}
    \end{center}
  \end{minipage}
\end{example}

\subsection{Coordonnées du vecteur $\wvec{AB}$}

\begin{proposition}{}{}
  \begin{multihide}
    Soit $A(x_A\,;\,y_A)$ et $B(x_B\,;\,y_B)$ deux points. Le vecteur $\wvec{AB}$ a pour coordonnées :
    \[
      \wvec{AB}\col{x_B - x_A}{y_B - y_A}
    \]
  \end{multihide}
\end{proposition}

\Needspace{8\baselineskip}
\begin{example}
  \begin{minipage}{0.55\linewidth}
    On donne $A(-2\,;\,3)$, $B(4\,;\,-1)$ et $C(-3\,;\,-2)$.
    \begin{itemize}
      \item $\wvec{AB}\col{x_B - x_A}{y_B - y_A} = \hideM{\col{4-(-2)}{-1-3}} = \hideM{\col{6}{-4}}$
      \item[]
      \item $\wvec{BC}\,\hideM{\col{-3-4}{-2-(-1)}} = \hideM{\col{-7}{-1}}$
      \item[]
      \item $\wvec{CA}\,\hideM{\col{-2-(-3)}{3-(-2)}} = \hideM{\col{1}{5}}$
    \end{itemize}
  \end{minipage}
  \hfill
  \begin{minipage}{0.42\linewidth}
    \begin{center}
    \begin{tikzpicture}[x=0.55cm, y=0.55cm]
      \axesG{-4}{5}{-3}{4}
      \pt{-2,3}{above left}{A} \pt{4,-1}{right}{B} \pt{-3,-2}{below left}{C}
      \hidetk{\draw[vec, blue!70!black] (-2,3) -- (4,-1);}
    \end{tikzpicture}
    \end{center}
  \end{minipage}
\end{example}

\begin{remark}
  Si $O$ est l'origine du repère, le vecteur $\wvec{OM}$ a les mêmes coordonnées que le point $M$. En physique, on l'appelle le \emph{vecteur position} de $M$, et $\wvec{AB}$ le \emph{vecteur déplacement} de $A$ à $B$.
\end{remark}

\subsection{Égalité de deux vecteurs}

\begin{proposition}{}{}
  \begin{multihide}
    Deux vecteurs sont égaux \ssi ils ont les \textbf{mêmes coordonnées}.
  \end{multihide}
\end{proposition}

\begin{method}{Trouver le quatrième sommet d'un parallélogramme}{}
  On rappelle que $ABCD$ est un parallélogramme \ssi $\wvec{AB} = \wvec{DC}$.

  On calcule les coordonnées de $\wvec{AB}$, puis on écrit que $\wvec{DC}$ a les mêmes coordonnées : on obtient une équation pour $x_D$ et une pour $y_D$.
\end{method}

\begin{example}
  On donne $A(1\,;\,1)$, $B(4\,;\,2)$ et $C(3\,;\,5)$. Déterminer les coordonnées du point $D$ tel que $ABCD$ soit un parallélogramme.
  \begin{reponse}
    $\wvec{AB}\col{4-1}{2-1} = \col{3}{1}$ et $\wvec{DC}\col{3 - x_D}{5 - y_D}$.

    $ABCD$ est un parallélogramme \ssi $\wvec{AB} = \wvec{DC}$, c'est-à-dire $3 - x_D = 3$ et $5 - y_D = 1$.

    Donc $x_D = 0$ et $y_D = 4$ : $D(0\,;\,4)$.
  \end{reponse}
\end{example}

\section{Distance entre deux points}

\begin{proposition}{Norme d'un vecteur}{}
  \begin{multihide}
    Dans un repère orthonormé, la norme du vecteur $\vec u\col{x}{y}$ est :
    \[
      \lVert \vec u \rVert = \sqrt{x^2 + y^2}
    \]
  \end{multihide}
\end{proposition}

\Needspace{8\baselineskip}
\begin{example}
  \begin{minipage}{0.55\linewidth}
    Pour $\vec u\col{4}{3}$, on trace un triangle rectangle dont les côtés de l'angle droit mesurent $4$ et $3$.

    D'après le théorème de Pythagore :
    \[
      \lVert \vec u \rVert^2 = \hideM{4^2 + 3^2 = 16 + 9 = 25}
    \]
    donc $\lVert \vec u \rVert = \hideM{\sqrt{25} = 5}$.
  \end{minipage}
  \hfill
  \begin{minipage}{0.4\linewidth}
    \begin{center}
    \begin{tikzpicture}[scale=0.7]
      \draw[gray!35, very thin, step=1] (-0.5,-0.5) grid (4.5,3.5);
      \draw[thick, dashed] (0,0) -- (4,0) node[midway, below] {$4$} -- (4,3) node[midway, right] {$3$};
      \draw (3.7,0) -- (3.7,0.3) -- (4,0.3);
      \draw[vec] (0,0) -- (4,3) node[midway, above left] {$\vec u$};
    \end{tikzpicture}
    \end{center}
  \end{minipage}
\end{example}

\begin{proposition}{Distance entre deux points}{}
  \begin{multihide}
    Dans un repère orthonormé, la distance entre $A(x_A\,;\,y_A)$ et $B(x_B\,;\,y_B)$ est la norme du vecteur $\wvec{AB}$ :
    \[
      AB = \lVert \wvec{AB} \rVert = \sqrt{(x_B - x_A)^2 + (y_B - y_A)^2}
    \]
  \end{multihide}
\end{proposition}

\begin{example}
  On reprend $A(-2\,;\,3)$, $B(4\,;\,-1)$ et $C(-3\,;\,-2)$, avec $\wvec{AB}\col{6}{-4}$, $\wvec{BC}\col{-7}{-1}$ et $\wvec{CA}\col{1}{5}$.
  \begin{itemize}
    \item $AB = \sqrt{6^2 + (-4)^2} = \hideM{\sqrt{36 + 16} = \sqrt{52} = 2\sqrt{13}}$
    \item $BC = \hideM{\sqrt{(-7)^2 + (-1)^2} = \sqrt{49 + 1} = \sqrt{50} = 5\sqrt{2}}$
    \item $CA = \hideM{\sqrt{1^2 + 5^2} = \sqrt{1 + 25} = \sqrt{26}}$
  \end{itemize}
\end{example}

\begin{method}{Déterminer la nature d'un triangle}{}
  \begin{itemize}
    \item On calcule les longueurs des trois côtés.
    \item Deux côtés de même longueur : le triangle est \textbf{isocèle}.
    \item Si le carré du plus grand côté est égal à la somme des carrés des deux autres, le triangle est \textbf{rectangle} (réciproque du théorème de Pythagore).
  \end{itemize}
\end{method}

\begin{example}
  On donne $E(1\,;\,1)$, $F(4\,;\,5)$ et $G(-3\,;\,4)$. Quelle est la nature du triangle $EFG$~?
  \begin{reponse}
    $EF = \sqrt{3^2 + 4^2} = \sqrt{25} = 5$ \quad et \quad $EG = \sqrt{(-4)^2 + 3^2} = \sqrt{25} = 5$.

    $FG = \sqrt{(-7)^2 + (-1)^2} = \sqrt{50}$.

    $EF = EG$ : le triangle est isocèle en $E$. De plus $EF^2 + EG^2 = 25 + 25 = 50 = FG^2$ : il est rectangle en $E$.

    Le triangle $EFG$ est donc rectangle isocèle en $E$.
  \end{reponse}
\end{example}

\section{Milieu d'un segment}

\begin{proposition}{Caractérisation vectorielle du milieu}{}
  \begin{multihide}
    $I$ est le milieu du segment $[AB]$ \ssi $\wvec{AI} = \wvec{IB}$.
  \end{multihide}
\end{proposition}

\begin{proposition}{Coordonnées du milieu}{}
  \begin{multihide}
    Le milieu $I$ du segment $[AB]$ a pour coordonnées :
    \[
      x_I = \frac{x_A + x_B}{2} \qquad \text{et} \qquad y_I = \frac{y_A + y_B}{2}
    \]
  \end{multihide}
\end{proposition}

\begin{demo}[Démonstration (pour aller plus loin)]
  Comme $\wvec{AI} = \wvec{IB}$, les deux vecteurs ont les mêmes coordonnées : $x_I - x_A = x_B - x_I$.
  Donc $2x_I = x_A + x_B$, soit $x_I = \dfrac{x_A + x_B}{2}$. On procède de même pour $y_I$.
\end{demo}

\begin{example}
  On reprend $A(-2\,;\,3)$, $B(4\,;\,-1)$ et $C(-3\,;\,-2)$. On note $I$, $J$ et $K$ les milieux de $[AB]$, $[BC]$ et $[AC]$.
  \begin{itemize}
    \item $x_I = \dfrac{-2 + 4}{2} = \hideM{1}$ \quad et \quad $y_I = \dfrac{3 + (-1)}{2} = \hideM{1}$, \quad donc $I\,\hideM{(1\,;\,1)}$.
    \item $J\,\hideM{\left(\dfrac{1}{2}\,;\,-\dfrac{3}{2}\right)}$ \qquad $K\,\hideM{\left(-\dfrac{5}{2}\,;\,\dfrac{1}{2}\right)}$
  \end{itemize}
\end{example}

\begin{method}{Utiliser le milieu}{}
  \begin{itemize}
    \item \textbf{Parallélogramme} : un quadrilatère est un parallélogramme \ssi ses diagonales ont le même milieu.
    \item \textbf{Symétrique} : $B$ est le symétrique de $A$ par rapport à $M$ \ssi $M$ est le milieu de $[AB]$.
  \end{itemize}
\end{method}

\begin{example}
  On donne $A(-3\,;\,1)$ et $M(2\,;\,-2)$. Déterminer les coordonnées du symétrique $B$ de $A$ par rapport à $M$.
  \begin{reponse}
    $M$ est le milieu de $[AB]$, donc $\dfrac{-3 + x_B}{2} = 2$ et $\dfrac{1 + y_B}{2} = -2$.

    D'où $-3 + x_B = 4$ et $1 + y_B = -4$, soit $x_B = 7$ et $y_B = -5$ : $B(7\,;\,-5)$.
  \end{reponse}
\end{example}

\section{Somme de vecteurs}

\begin{proposition}{}{}
  \begin{multihide}
    Si $\vec u\col{x}{y}$ et $\vec v\col{x'}{y'}$, alors $\vec u + \vec v\col{x + x'}{y + y'}$.
  \end{multihide}
\end{proposition}

\Needspace{8\baselineskip}
\begin{example}
  \begin{minipage}{0.5\linewidth}
    On a représenté deux vecteurs $\vec u$ et $\vec v$.
    \begin{itemize}
      \item $\vec u\,\hideM{\col{4}{1}}$ \quad et \quad $\vec v\,\hideM{\col{-1}{-3}}$
      \item[]
      \item $\vec u + \vec v\,\hideM{\col{4 + (-1)}{1 + (-3)} = \col{3}{-2}}$
    \end{itemize}
    On construit $\vec u + \vec v$ en mettant les flèches bout à bout : on retrouve les mêmes coordonnées.
  \end{minipage}
  \hfill
  \begin{minipage}{0.45\linewidth}
    \begin{center}
    \begin{tikzpicture}[x=0.6cm, y=0.6cm]
      \draw[gray!35, very thin, step=1] (-1,-1) grid (7,4);
      \draw[vec] (0,2) -- (4,3) node[midway, above] {$\vec u$};
      \draw[vec] (6,3) -- (5,0) node[midway, right] {$\vec v$};
      \hidetk{
        \draw[vec, dashed] (4,3) -- (3,0);
        \draw[vec, blue!70!black] (0,2) -- (3,0) node[midway, below left] {$\vec u + \vec v$};
      }
    \end{tikzpicture}
    \end{center}
  \end{minipage}
\end{example}

\begin{remark}
  On retrouve la relation de Chasles. Avec $A(-2\,;\,3)$, $B(4\,;\,-1)$ et $C(-3\,;\,-2)$ :
  \[
    \wvec{AB} + \wvec{BC} = \col{6}{-4} + \col{-7}{-1} = \hideM{\col{-1}{-5}}
    \qquad \text{et} \qquad
    \wvec{AC}\,\hideM{\col{-3-(-2)}{-2-3} = \col{-1}{-5}}
  \]
\end{remark}

\section{Produit d'un vecteur par un réel}

\begin{definition}{Rappel du chapitre 1}{}
  \begin{multihide}
    Soit $\vec u$ un vecteur non nul et $k$ un réel non nul. Le vecteur $k\vec u$ a :
    \begin{itemize}
      \item la \textbf{même direction} que $\vec u$ ;
      \item le \textbf{même sens} que $\vec u$ si $k > 0$, le \textbf{sens contraire} si $k < 0$ ;
      \item une norme égale à $\lvert k \rvert \times \lVert \vec u \rVert$.
    \end{itemize}
  \end{multihide}
\end{definition}

\begin{example}
  À partir du vecteur $\vec u$, construire les vecteurs $2\vec u$, $-\vec u$ et $\dfrac{1}{2}\vec u$.

  \begin{center}
  \begin{tikzpicture}[x=0.6cm, y=0.6cm]
    \draw[gray!35, very thin, step=1] (0,0) grid (22,5);
    \draw[vec] (1,1) -- (5,3) node[midway, above left] {$\vec u$};
    \hidetk{
      \draw[vec, blue!70!black] (6,1) -- (14,5) node[midway, above left] {$2\vec u$};
      \draw[vec, blue!70!black] (19,3) -- (15,1) node[midway, above left] {$-\vec u$};
      \draw[vec, blue!70!black] (19,4) -- (21,5) node[midway, above left] {$\tfrac12\vec u$};
    }
  \end{tikzpicture}
  \end{center}
\end{example}

\begin{proposition}{}{}
  \begin{multihide}
    Si $\vec u\col{x}{y}$ et $k$ est un réel, alors $k\vec u\col{kx}{ky}$.
  \end{multihide}
\end{proposition}

\begin{example}
  Avec $\vec u\col{4}{-2}$ et $\vec v\col{-1}{3}$ :
  \begin{tasks}[label=\textbullet](2)
    \task $3\vec u\,\hideM{\col{3 \times 4}{3 \times (-2)} = \col{12}{-6}}$
    \task $-\dfrac{1}{2}\vec u\,\hideM{\col{-2}{1}}$
    \task $2\vec u - 3\vec v\,\hideM{\col{8 - (-3)}{-4 - 9} = \col{11}{-13}}$
    \task $\vec u + 4\vec v\,\hideM{\col{4 - 4}{-2 + 12} = \col{0}{10}}$
  \end{tasks}
\end{example}

\begin{example}
  On donne $A(1\,;\,2)$ et $B(2\,;\,-1)$. Déterminer les coordonnées du point $M$ tel que $\wvec{AM} = 3\wvec{AB}$.
  \begin{reponse}
    $\wvec{AB}\col{1}{-3}$, donc $3\wvec{AB}\col{3}{-9}$. Et $\wvec{AM}\col{x_M - 1}{y_M - 2}$.

    $x_M - 1 = 3$ et $y_M - 2 = -9$, donc $M(4\,;\,-7)$.
  \end{reponse}
\end{example}

\vfill

\begin{savoir}{}{}
\begin{itemize}[label=--]
  \item Lire les coordonnées d'un vecteur et représenter un vecteur de coordonnées données.
  \item Calculer les coordonnées d'un vecteur $\wvec{AB}$.
  \item Déterminer le quatrième sommet d'un parallélogramme.
  \item Calculer une distance et déterminer la nature d'un triangle.
  \item Calculer les coordonnées du milieu d'un segment, d'un symétrique.
  \item Calculer les coordonnées d'une somme de vecteurs et d'un produit par un réel.
  \item Construire le vecteur $k\vec u$.
\end{itemize}
\end{savoir}

\end{document}
